\subsection{What Quorum Is}

Quorum is an AI-powered GitHub Pull Request review system. It is a web
application that automatically analyzes pull requests opened in a GitHub
repository, produces a deterministic Merge Readiness Score, and presents the
results in a web dashboard with an evidence-grounded question-answering
assistant called \emph{Ask Quorum}. The system is implemented as a Python /
FastAPI backend, a React/TypeScript frontend, and a PostgreSQL database, and
it integrates with GitHub through a GitHub App, OAuth authentication, and
webhooks.

\subsection{The Problem Addressed}

Reviewing a pull request manually is time-consuming and error-prone. A
reviewer must inspect the changed code, look for security problems, estimate
whether the change might break existing behaviour, and form an opinion about
whether the change is safe to merge. This judgement is often based on
subjective impressions rather than measured evidence, and it is easy to miss
a security issue or a failing test. The conventional review process also
requires the reviewer to either write and run tests manually or trust the
author's claims.

\subsection{What Quorum Automates}

Quorum automates the evidence-gathering part of code review. When a pull
request event reaches the backend, the orchestrator runs a review pipeline
that:

\begin{itemize}
  \item extracts and parses the pull request diff and the changed Python
        files (diff and AST analysis);
  \item runs the \textbf{Semgrep} static analysis tool over the changed files
        to obtain security evidence;
  \item passes the Semgrep findings and code context to a \textbf{Security
        Agent} (an LLM-driven agent) which reasons over the evidence and
        produces structured, evidence-traceable security findings;
  \item uses a \textbf{Test Writer Agent} to generate \texttt{pytest} tests
        for the modified functions;
  \item executes the generated tests inside an isolated, hardened
        \textbf{Docker sandbox} and measures code coverage before and after;
  \item combines the security, test, and coverage results into a
        deterministic \textbf{Merge Readiness Score} from 0 to 100;
  \item persists the results in PostgreSQL and posts a concise review summary
        back to the pull request on GitHub.
\end{itemize}

\subsection{How the Result Helps Developers}

The Merge Readiness Score is deterministic: the same evidence always produces
the same score, and the breakdown (security findings, generated tests, and
coverage) is stored with the analysis. A developer can therefore look at a
pull request and immediately understand both the score and the reasons behind
it, rather than relying on an unverifiable subjective opinion. The dashboard
shows security findings with evidence, generated test outcomes, coverage
deltas, and the score for each pull request, and the Ask Quorum assistant
answers questions about a specific review using only that review's stored
evidence.

\subsection{The Role of the Dashboard}

The React dashboard is the visual control centre for the pipeline. It allows a
signed-in user to discover and connect GitHub repositories, open a repository
workspace, create and manage pull requests (including creating, merging, and
closing them from Quorum), view per-pull-request analysis results, review
history, security findings, and settings, and ask questions about a review.

\subsection{Project Objective}

The overall objective of the project is to provide an end-to-end,
evidence-based pull request review tool that helps developers assess whether a
change is safe to merge. The project demonstrates the combination of static
security analysis, LLM-driven reasoning, automated test generation, sandboxed
test execution, coverage measurement, and deterministic result synthesis, all
presented through a multi-user web application integrated with GitHub.

\subsection{Implementation Status}

All planned roadmap phases (Phase 0 through Phase 19) are implemented. The
LLM layer is active through the OpenAI provider; a local \texttt{Ollama}
provider exists in the codebase but is not currently used. Components that are
planned but not yet implemented (such as an automated CI pipeline and a backup
demo video) are noted in the project documentation as future work.